\documentclass[11pt]{amsart}
\usepackage{geometry}                % See geometry.pdf to learn the layout options. There are lots.
\geometry{letterpaper}                   % ... or a4paper or a5paper or ... 
%\geometry{landscape}                % Activate for for rotated page geometry
%\usepackage[parfill]{parskip}    % Activate to begin paragraphs with an empty line rather than an indent
\usepackage{graphicx}
\usepackage{amssymb}
\usepackage{epstopdf}
\usepackage{multirow}
\DeclareGraphicsRule{.tif}{png}{.png}{`convert #1 `dirname #1`/`basename #1 .tif`.png}

\title{Problema PS 2.1}
\author{}
%\date{}                                           % Activate to display a given date or no date

\begin{document}
\maketitle
%\section{}
%\subsection{}
La tangente de un ángulo se define como el cociente entre el seno y el coseno de dicho ángulo:

\[
\tan(a) = \frac{\sin(a)}{\cos(a)}
\]

Construya un diagrama de flujo que le permita calcular la tangente de un ángulo, considerando que se conoce el valor del seno y del coseno del mismo. Recuerde que el coseno debe ser diferente de 0.

Datos: SENO, COSENO (variables de tipo real).

\end{document}


